% Objective — Data in Brief body fragment.
% Section header (\section{Objective}) is supplied by the master dib.tex.
% Descriptive only: motivates the dataset, names the parent JSR article, cites
% RASAero II and OpenRocket. NO new scientific conclusions, NO RASAero result
% argument (that belongs to the companion JSR paper). All numbers match
% paper/_shared/CANONICAL_FACTS.md and the released CSVs.

Validating a flight-trajectory simulator in the supersonic regime requires a
ground-truth corpus that pairs measured flight outcomes with the predictions of
an established reference tool. For hobby- and sounding-rocket aerodynamics, no
such corpus has been openly available: published flight records are scattered
across primary sources in differing formats, and the most widely used commercial
predictions for these vehicles---those produced by RASAero~II~\cite{rogers2015}---are
not distributed as a machine-readable, reusable table alongside the corresponding
measured apogees. As a result, comparisons between open-source simulators such as
OpenRocket~\cite{niskanen2009} and a commercial baseline have typically been made
on small, internally selected sets of flights, which makes the resulting accuracy
statistics difficult to reproduce and easy to bias through outcome-based selection.

The Rocket Flight Database addresses this gap. Its headline component is a
25-flight corpus spanning Mach 0.54--4.33, drawn from Rogers' publicly published
RASAero~II altitude-comparison collection (23 single-stage high-power flights plus
two two-stage flights---the AeroPac~104K Two-Stage at Mach~3.04 and the MESOS~293K
flight, the highest-Mach case at Mach~4.33). Because the
flights were selected externally---by the authors of the commercial comparison set
rather than curated by us for favourable outcomes---and because each record carries
both the measured apogee and the paired RASAero~II prediction, the table supports an
outcome-independent, reproducible validation of supersonic apogee prediction. The
ground-truth column reports instrumented apogees (barometric, GPS, optical-track,
and integrated-accelerometer records), and the commercial-reference column reproduces
Rogers' recorded RASAero~II predictions verbatim; only the OpenRocket-Plus prediction
column is regenerated from the archived simulator code, so the comparison is
non-circular.

A companion exploratory table extends the database to roughly twenty historical
high-Mach sounding-rocket flights (Black Brant~V, Nike--Apache, Nike--Cajun,
Nike--Deacon, Arcas, and HEROS~3), reported in full---including the cases that fall
well outside the headline admission band---to document the simulator's behaviour
beyond Mach~4.33 without selection bias.

This dataset is the validation corpus underlying the companion research article that
introduces OpenRocket-Plus, an open-source supersonic extension of
OpenRocket~\cite{niskanen2009}, and reports its integrated apogee accuracy against
the RASAero~II baseline~\cite{yu2026jsr}. That parent article presents and interprets
the comparison result; the present article is descriptive and confines itself to the
data's schema, provenance, and reuse value. Releasing the corpus openly under a
permissive licence lets other investigators independently reproduce those statistics,
benchmark alternative tools---including newer open-source trajectory frameworks such
as RocketPy~\cite{rocketpy2021} and recent amateur-supersonic flight studies
\cite{quintart2025,lowcostroll2025}---and extend the corpus with additional
instrumented flights.
